%Martyna Boniatowska
\documentclass[11pt]{article}
\usepackage[utf8]{inputenc}
\usepackage{polski}
\usepackage[sf]{titlesec}
\usepackage{graphicx}
\usepackage{geometry}
\usepackage{fancybox}
\usepackage{enumitem}
\usepackage{ifpdf}
\usepackage{fontenc}
\usepackage{url}
\usepackage{amsmath}
\usepackage{mathtools}
\geometry{
 a4paper,
 total={170mm,257mm},
 left=20mm,
 top=20mm,
 }
\usepackage[dvipsnames]{xcolor}%niestandardowe kolory tekstu
\usepackage{fontawesome5}%ikonki

\newcommand*{\myfont}{\fontfamily{qag}\selectfont}
\newenvironment{my}{\fontfamily{pbk}\selectfont}{\par}
\titleformat{\section}         % Customise the \section command 
  {\color{BrickRed}\Large\myfont\raggedright} % Make the \section headers large 
                               % small capitals (\scshape) and left aligned (\raggedright)
  {}{0em}                      % Can be used to give a prefix to all sections, like 'Section ...'
  {\MakeUppercase}                           % Can be used to insert code before the heading
  [\color{Gray}{\titlerule[1pt]}]                 % Inserts a horizontal line after the heading

\titleformat{\subsection}
  {\normalfont\scshape\raggedright}
  {}{0em}
  {}

\newcommand{\datedsubsection}[2]{%
  \subsection[#1]{#1 \hfill \normalfont{#2}}%
}
\newcommand{\kropki}[5]{\Large #1 $\bullet$ #2 $\bullet$ #3 $\bullet$ #4 $\bullet$ #5 $\bullet$}
\newcommand{\br}{\color{BrickRed}}
\newcommand{\g}{\color{gray}}
\begin{document}
\thispagestyle{empty} 
\begin{minipage}{0.25\textwidth}
\ifpdf
\includegraphics[scale = 0.02]{zdj1.jpg}
\else
\includegraphics[scale = 0.02]{zdj_cv.eps}
\fi
\end{minipage}
\begin{minipage}{0.75\textwidth}
\vskip 15pt
{\color{BrickRed}\huge\myfont Martyna Maria Boniatowska}\\[15pt]
{\color{BrickRed}\large\myfont Absolwentka Matematyki Stosowanej}\\[15pt]
\faIcon{home} :  \texttt{Wrocław} \\[10pt]
\faIcon{phone-square-alt} :  \texttt{572 454 821}\\[10pt]
\faIcon{envelope} : \texttt{mboniatowska380@gmail.com}\\[10pt]
\faIcon{github} : \url{github.com/mboniatowska/Projects}
\end{minipage}
\newline
\begin{my}
\large{
\section{Edukacja}
\datedsubsection{Politechnika Wrocławska}{10.2018 -- 02.2022}
Matematyka Stosowana -- studia inżynierskie.
\begin{itemize}[label={}]
    \item {\color{BrickRed}$\triangleright$} Temat pracy dyplomowej: Analiza korelacji, regresja liniowa i regresja logistyczna w badaniach sportowych.
\end{itemize}
\datedsubsection{Politechnika Wrocławska}{02.2022 -- 06.2023}
Applied Mathematics with specialization in Data Engineering -- studia magisterskie.
%%%%%%%%%%%%%%%%%%%%%%%%%%%%%%%%%%%%%%%%%%%%%%%%%%%%%%%%%%%%%%%%%%%%%
\section{Doświadczenie}
\datedsubsection{Marbet Lipsko}{03.2021 --- 04.2021}
Praktyka studencka w dziale informatyki i analizy danych.
\begin{itemize}[label={}]
    \item {\color{BrickRed}$\triangleright$} Wykorzystanie SQL'a do tworzenia raportów i przeprowadzania analiz.
    \item {\color{BrickRed}$\triangleright$} Porządkowanie dokumentacji.

\end{itemize}
%%%%%%%%%%%%%%%%%%%%%%%%%%%%%%%%%%%%%%%%%%%%%%%%%%%%%%%%%%%%%%%%%%%%%
\section{Oprogramowanie}
\begin{tabular}{p{2.5cm} c p{4cm} p{2.5cm} c}
Python & \kropki{\br}{\br}{\br}{\g}{\g} &  &
LaTeX & \kropki{\br}{\br}{\br}{\br}{\g}\\
Matlab & \kropki{\br}{\br}{\br}{\br}{\g} &  & 
Excel &\kropki{\br}{\br}{\br}{\g}{\g}\\
SQL  &\kropki{\br}{\br}{\br}{\g}{\g}   & &
MS Office &\kropki{\br}{\br}{\br}{\br}{\g}
\end{tabular}
%%%%%%%%%%%%%%%%%%%%%%%%%%%%%%%%%%%%%%%%%%%%%%%%%%%%%%%%%%%%%%%%%%%
\section{Umiejętności miękkie}
{\ovalbox{\color{BrickRed} Wizualizacja danych}  \ovalbox{\color{BrickRed} Umiejętność prezentacji} \ovalbox{\color{BrickRed} Zdolność rozwiązywania problemów}\\[5pt] \ovalbox{\color{BrickRed} Zarządzanie czasem} \ovalbox{\color{BrickRed} Kreatywność} \ovalbox{\color{BrickRed} Zdolności interpersonalne} \ovalbox{\color{BrickRed} Odpowiedzialność}\\[5pt]
\ovalbox{\color{BrickRed} Dbałość o szczegóły} \ovalbox{\color{BrickRed} Praca zespołowa}}

%%%%%%%%%%%%%%%%%%%%%%%%%%%%%%%%%%%%%%%%%%%%%%%%%%%%%%%%%%%%%%%%%%%
\section{Języki obce}
\begin{tabular}{p{2.5cm} c p{4cm} p{2.5cm} c}
Angielski & \kropki{\br}{\br}{\br}{\br}{\g} &  &
\end{tabular}
%%%%%%%%%%%%%%%%%%%%%%%%%%%%%%%%%%%%%%%%%%%%%%%%%%%%%%%%%%%%%%%%%%%%
}
\end{my}

\end{document}